A public vulnerability fix teaches a lifecycle rule: an object must not
be released without resetting a caller-visible reference to it, or
without detaching it from a chain it still belongs to. Existing
general-purpose static analyzers do not check a specific project's own
such conventions; they check their own stock rule sets instead.
CVE-2022-4450 is a real case: \texttt{PEM\_read\_bio\_ex} in OpenSSL frees
two output buffers on one failure path without clearing the
caller-visible pointers to them, letting a caller double-free. We present
\tool, which reads such a rule from a data file and checks it against a
\texttt{tree-sitter} control-flow graph. On OpenSSL, \tool flags \texttt{PEM\_read\_bio\_ex}'s two outputs
before the fix and clears both after it. It also reproduces an unrelated
non-CVE fix the rule was never shown, and flags one site a later OpenSSL
commit went on to fix independently.
On the current 1{,}155-file \texttt{crypto/ssl/providers} tree it flags
10 sites: manual triage finds 2 potential recurrences whose failure
branch is not reachable in normal builds, 8 harmless deviations, and 0
confirmed new vulnerabilities. None of three
general-purpose open-source static analyzers we ran, with their stock
configuration, distinguishes any of three historical CVEs from their own
fix; a project-specific rule, not a sharper general one, closed that gap.
A second, differently-shaped rule for the detach-before-release BIO-chain
invariant behind CVE-2023-0215 reuses the same parser and control-flow
graph and also distinguishes its own motivating fix. Its whole-tree run,
though, flags 11 sites that are all false positives, because the rule has
no notion of which released object is externally owned. 
